\documentclass[12pt,letterpaper]{article}

% ── Packages ──────────────────────────────────────────────────────────────────
\usepackage[utf8]{inputenc}
\usepackage[T1]{fontenc}
\usepackage[english]{babel}
\usepackage{amsmath,amssymb,amsthm}
\usepackage{mathtools}
\usepackage{geometry}
\usepackage{graphicx}
\usepackage{xcolor}
\usepackage{hyperref}
\usepackage{booktabs}
\usepackage{array}
\usepackage{listings}
\usepackage{algorithm}
\usepackage{algpseudocode}
\usepackage{tikz}
\usepackage{cite}
\usepackage{caption}
\usepackage{subcaption}
\usepackage{float}
\usepackage{enumitem}
\usepackage{fancyhdr}
\usepackage{setspace}

\usetikzlibrary{matrix,arrows.meta,positioning,shapes,calc}

% ── Geometry ──────────────────────────────────────────────────────────────────
\geometry{
    headheight=14pt,
    letterpaper,
    top=2.5cm, bottom=2.5cm,
    left=3cm,  right=2.5cm
}

% ── Colors ────────────────────────────────────────────────────────────────────
\definecolor{azulnodo}{RGB}{33,150,243}
\definecolor{verdenodo}{RGB}{76,175,80}
\definecolor{naranjanodo}{RGB}{255,152,0}
\definecolor{rosanodo}{RGB}{233,30,99}
\definecolor{grisoscuro}{RGB}{55,71,79}
\definecolor{codeblue}{RGB}{0,100,180}
\definecolor{codegray}{RGB}{120,120,120}

% ── Listings ──────────────────────────────────────────────────────────────────
\lstset{
    language=Python,
    basicstyle=\ttfamily\small,
    keywordstyle=\color{codeblue}\bfseries,
    commentstyle=\color{codegray}\itshape,
    stringstyle=\color{verdenodo},
    numbers=left,
    numberstyle=\tiny\color{codegray},
    stepnumber=1,
    numbersep=8pt,
    backgroundcolor=\color{gray!8},
    showstringspaces=false,
    breaklines=true,
    frame=single,
    framerule=0.4pt,
    rulecolor=\color{gray!40},
    captionpos=b,
    tabsize=4,
}

% ── Theorem environments ───────────────────────────────────────────────────────
\theoremstyle{plain}
\newtheorem{theorem}{Theorem}[section]
\newtheorem{lemma}[theorem]{Lemma}
\newtheorem{corollary}[theorem]{Corollary}

\theoremstyle{definition}
\newtheorem{definition}[theorem]{Definition}

\theoremstyle{remark}
\newtheorem{remark}[theorem]{Remark}

% ── Header/Footer ─────────────────────────────────────────────────────────────
\pagestyle{fancy}
\fancyhf{}
\fancyhead[L]{\small\textit{Planar Graph Coloring --- Topological Peeling and CSP}}
\fancyhead[R]{\small\thepage}
\fancyfoot[C]{\small Programa Delf\'in 2026 --- UAEM}
\renewcommand{\headrulewidth}{0.4pt}
\renewcommand{\footrulewidth}{0.4pt}

% ── Hyperref ──────────────────────────────────────────────────────────────────
\hypersetup{
    colorlinks=true,
    linkcolor=codeblue,
    citecolor=grisoscuro,
    urlcolor=codeblue,
    pdftitle={Planar Graph Coloring via Topological Peeling and CSP Backtracking},
    pdfauthor={Yosef Ali Jimenez Munoz}
}

% ── Misc ──────────────────────────────────────────────────────────────────────
\setlength{\parskip}{6pt}
\setlength{\parindent}{1.2em}

\newcommand{\N}{\mathbb{N}}
\newcommand{\Z}{\mathbb{Z}}
\newcommand{\R}{\mathbb{R}}
\newcommand{\degr}{\operatorname{deg}}
\newcommand{\ecc}{\operatorname{ecc}}
\newcommand{\dom}{\operatorname{dom}}
\newcommand{\argmin}{\operatorname{argmin}}

% ══════════════════════════════════════════════════════════════════════════════
\begin{document}
% ══════════════════════════════════════════════════════════════════════════════

% ── Title Page ────────────────────────────────────────────────────────────────
\begin{titlepage}
\centering
\vspace*{2cm}

{\large Universidad Aut\'onoma del Estado de M\'exico}\\[0.3cm]
{\large Programa Delf\'in 2026}\\[0.3cm]
{\large Research Line: Theory of Computation}\\[2cm]

{\LARGE\bfseries Planar Graph Coloring}\\[0.4cm]
{\LARGE\bfseries via Topological Peeling}\\[0.4cm]
{\LARGE\bfseries and CSP Backtracking}\\[0.8cm]

{\large Mathematical and algorithmic documentation\\
of a computational visualization of\\
the Four Color Theorem}\\[2.5cm]

\begin{tabular}{ll}
\textbf{Author:}              & Yosef Ali Jim\'enez Mu\~noz \\
\textbf{Student ID:}          & 492410631 \\
\textbf{Home Institution:}    & Universidad Tecnol\'ogica de Altamira \\
\textbf{Academic Advisor:}    & Dr. Jos\'e Raymundo Marcial Romero \\
\textbf{Program:}             & Programa Delf\'in, Summer 2026 \\
\textbf{Date:}                & June 2026
\end{tabular}

\vfill
\end{titlepage}

% ── Abstract ──────────────────────────────────────────────────────────────────
\begin{abstract}
We document the mathematical and algorithmic design of a computational
visualization system for the Four Color Theorem, implemented in Python with the
libraries \texttt{NetworkX}, \texttt{NumPy}, \texttt{SciPy}, and
\texttt{Matplotlib}. The system operates on arbitrary planar graphs and executes
a transformation pipeline that includes: planarity verification via the
Boyer-Myrvold algorithm~\cite{boyermyrvold1999}; straight-line drawing via the
Tutte embedding~\cite{tutte1963}; topological layer decomposition via outer-face
peeling; inverse-order greedy coloring; analysis of Kempe
chains~\cite{kempe1879}; and CSP backtracking with MRV and forward checking when
the greedy fails to produce four colors. The Errera~\cite{errera1921} and
Kittell~\cite{kittell1935} graphs are used as canonical test cases, concretely
demonstrating the failure mechanism of Kempe's algorithm via tangled chains. The
document includes complete mathematical foundations for each step, precise
numerical values computed on the Errera graph, and situates the work historically
from Kempe's incorrect proof to Gonthier's formal verification in
Coq~\cite{gonthier2008}.
\end{abstract}

\tableofcontents
\newpage

% ══════════════════════════════════════════════════════════════════════════════
\section{Introduction}
% ══════════════════════════════════════════════════════════════════════════════

The Four Color Theorem states that every planar graph can be colored with at
most four colors so that no two adjacent vertices share a color. Formulated as a
conjecture in 1852 by Francis Guthrie while coloring a map of English counties,
it resisted proof for over a century~\cite{wilson2002}.

The first proof was proposed by Alfred Kempe in 1879~\cite{kempe1879}, but Percy
Heawood identified a fundamental error in 1890~\cite{heawood1890}. In 1921,
Alfred Errera constructed a concrete planar graph exhibiting the exact failure
point of Kempe's algorithm~\cite{errera1921}. A correct proof did not arrive
until 1976, when Appel and Haken computationally verified that a finite set of
configurations covers all possible cases~\cite{appelhaken1977}. Robertson,
Sanders, Seymour, and Thomas simplified the proof to 633 configurations in
1997~\cite{robertson1997}, and in 2005 Gonthier produced the first mechanically
checkable formal verification in the Coq proof
assistant~\cite{gonthier2008}.

This document describes a visualization system that implements the algorithmic
steps of the coloring process, from the representation of a graph as ordered
pairs to the guarantee of four colors via CSP backtracking. The goal is not to
prove the theorem (whose proof requires the global discharging argument) but to
make visible the mathematical mechanisms underlying the problem: Euler's formula,
triangulation, barycentric embedding, topological peeling, greedy coloring,
Kempe chain failure, and constraint-based search.

% ══════════════════════════════════════════════════════════════════════════════
\section{Graph Representation}
% ══════════════════════════════════════════════════════════════════════════════

\subsection{Formal Definition}

\begin{definition}[Simple undirected graph]
A graph $G = (V, E)$ consists of a finite set of vertices $V$ and a set of edges
$E \subseteq \binom{V}{2}$, where $\binom{V}{2}$ denotes all two-element subsets
of $V$. Each edge $e = \{u,v\} \in E$ is an unordered pair with $u \neq v$.
\end{definition}

The \textbf{degree} of a vertex $v$ is $\degr(v) = |\{u \in V : \{u,v\} \in E\}|$.
The following identity holds because each edge contributes 1 to the degree of
each of its two endpoints:

\begin{equation}
    \sum_{v \in V} \degr(v) = 2|E|.
    \label{eq:handshake}
\end{equation}

\subsection{Input Format}

The system reads a \texttt{grafo.json} file with the structure:

\begin{lstlisting}[caption={Input file structure},label={lst:json}]
{
  "vertices": {"0": [x0, y0], "1": [x1, y1], ...},
  "aristas":  [[0, 1], [1, 2], ...]
}
\end{lstlisting}

The coordinates $(x_i, y_i)$ are used initially for geometric planarity checks,
but the system then computes a canonical embedding via the Tutte algorithm
(Section~\ref{sec:tutte}). Conversion from the edge list to the \texttt{NetworkX}
adjacency structure runs in $O(|E|)$.

\subsection{The Errera Graph}

The Errera graph~\cite{errera1921} has $|V| = 17$ vertices and $|E| = 45$
edges. Its degree sequence is twelve vertices of degree 5 and five of degree 6.
Verification: $12 \cdot 5 + 5 \cdot 6 = 90 = 2 \cdot 45$,
consistent with~\eqref{eq:handshake}. Its vertex connectivity is
$\kappa(G) = 5$~\cite{errera1921}.

The graph has a concentric layered structure:

\begin{center}
\begin{tabular}{clc}
\toprule
\textbf{Layer} & \textbf{Vertices} & \textbf{Role} \\
\midrule
Bottom apex & $\{0\}$            & Degree-5 vertex \\
Ring 1      & $\{1,2,3,4,5\}$   & Inner pentagon \\
Ring 2      & $\{6,7,8,9,10\}$  & Middle pentagon \\
Ring 3      & $\{11,12,13,14,15\}$ & Outer pentagon \\
Top apex    & $\{16\}$           & Degree-5 vertex \\
\bottomrule
\end{tabular}
\end{center}

This structure corresponds to two stacked pentagonal antiprisms with a
5-fold dihedral symmetry group $D_5$~\cite{errera1921}.

% ══════════════════════════════════════════════════════════════════════════════
\section{Planarity Verification}
% ══════════════════════════════════════════════════════════════════════════════

\subsection{Definition and Fundamental Theorems}

\begin{definition}[Planar graph]
A graph $G$ is planar if there exists an embedding $\varphi: G \to \mathbb{R}^2$
such that each vertex maps to a distinct point, each edge maps to a Jordan curve
between its endpoints, and curves of distinct edges do not intersect except
possibly at shared endpoints.
\end{definition}

\begin{theorem}[F\'ary, 1948~\cite{fary1948}]
Every simple planar graph admits a straight-line embedding: a drawing where
every edge is a line segment.
\end{theorem}

F\'ary's theorem justifies working with coordinates and segments without loss of
generality.

\begin{theorem}[Kuratowski, 1930~\cite{kuratowski1930}]
A graph is planar if and only if it contains no subgraph homeomorphic to $K_5$
or $K_{3,3}$.
\end{theorem}

\subsection{Boyer-Myrvold Algorithm}

The system calls \texttt{nx.check\_planarity(G)}, which internally runs the
Boyer-Myrvold algorithm~\cite{boyermyrvold1999}.

\begin{theorem}[Boyer-Myrvold, 1999~\cite{boyermyrvold1999}]
The Boyer-Myrvold algorithm determines planarity in $O(|V|)$ time and, when the
graph is planar, constructs the combinatorial embedding.
\end{theorem}

The algorithm operates on the combinatorial representation, not on coordinates.
For each vertex $v$ it stores the cyclic order $(\sigma_v(e_1), \sigma_v(e_2),
\ldots)$ of incident edges, where $\sigma_v$ is a cyclic permutation. This
structure---the \textbf{combinatorial embedding}---is sufficient to determine
faces without coordinates.

Geometric cross-checking with \texttt{Shapely} is an additional validation step
applied to the input JSON. The intersection test between two segments
$\overline{p_1 p_2}$ and $\overline{p_3 p_4}$ is based on the sign of the
cross product:

\begin{equation}
(p_2 - p_1) \times (p_3 - p_1) \quad \text{and} \quad (p_2 - p_1) \times (p_4 - p_1).
\end{equation}

The segments intersect in their interiors if and only if the signs differ in both
reciprocal checks~\cite{shamos1976}.

% ══════════════════════════════════════════════════════════════════════════════
\section{Euler's Formula and Maximality}
% ══════════════════════════════════════════════════════════════════════════════

\subsection{Euler's Formula}

\begin{theorem}[Euler, 1752~\cite{euler1752}]
For every connected planar graph, letting $F$ denote the number of faces
(regions, including the unbounded exterior face):
\begin{equation}
    V - E + F = 2.
    \label{eq:euler}
\end{equation}
\end{theorem}

\begin{proof}[Proof sketch]
By induction on $|E|$. A tree with $n$ vertices has $n-1$ edges and one face,
giving $n - (n-1) + 1 = 2$. Adding an edge between two components:
$\Delta E = 1$, $\Delta F = 0$, so $V - E + F$ is preserved. Adding an edge
within a component: $\Delta E = 1$, $\Delta F = 1$, also preserved.
\end{proof}

For the Errera graph: $17 - 45 + 30 = 2$ $\checkmark$, with 30 triangular faces.

\subsection{Density Bound}

\begin{theorem}
\label{thm:density}
For every simple planar graph with $|V| \geq 3$:
\begin{equation}
    |E| \leq 3|V| - 6.
    \label{eq:density}
\end{equation}
\end{theorem}

\begin{proof}
Every face has at least 3 edges on its boundary, and each edge borders exactly 2
faces, so $3F \leq 2|E|$, i.e., $F \leq \frac{2|E|}{3}$.
Substituting into~\eqref{eq:euler}:
\[
    |V| - |E| + \frac{2|E|}{3} \geq 2
    \implies |V| - \frac{|E|}{3} \geq 2
    \implies |E| \leq 3|V| - 6. \qedhere
\]
\end{proof}

\subsection{Maximal Planar Graphs}

\begin{definition}[Maximal planar graph]
A planar graph $G$ is maximal if for every $\{u,v\} \notin E$, the graph
$G' = (V, E \cup \{\{u,v\}\})$ is non-planar.
\end{definition}

\begin{theorem}
A planar graph with $|V| \geq 3$ is maximal if and only if every face (including
the outer face) is a triangle.
\end{theorem}

When the outer face is also a triangle, equality holds in~\eqref{eq:density}:
$|E| = 3|V| - 6$. In the general case where the outer face has $h$ vertices:
\begin{equation}
    |E| = 3|V| - 6 - (h - 3).
    \label{eq:maximal_general}
\end{equation}

For the Errera graph: $|E| = 45 = 3(17) - 6 = 45$ $\checkmark$, with $h = 3$.
The graph is already maximal, so the maximalization step adds zero edges.

\subsection{Existence of a Low-Degree Vertex}

\begin{theorem}
\label{thm:deg5}
In every planar graph there exists at least one vertex $v$ with $\degr(v) \leq 5$.
\end{theorem}

\begin{proof}
Suppose $\degr(v) \geq 6$ for all $v$. Then:
\[
2|E| = \sum_{v} \degr(v) \geq 6|V| \implies |E| \geq 3|V|.
\]
But Theorem~\ref{thm:density} requires $|E| \leq 3|V| - 6 < 3|V|$. Contradiction.
\end{proof}

This result guarantees that the inverse greedy coloring always finds a free color
among six (Section~\ref{sec:greedy}).

\subsection{Constrained Maximalization}

The input graph may not be maximal. The system adds edges without crossing any
existing ones:

\begin{algorithm}[H]
\caption{Constrained maximalization}
\label{alg:maximal}
\begin{algorithmic}[1]
\Require Planar graph $G = (V, E)$ with positions $\mathbf{p}: V \to \mathbb{R}^2$
\Ensure Maximal planar graph $G'$ with $E \subseteq E'$
\State $\mathcal{C} \gets \binom{V}{2} \setminus E$
\State Sort $\mathcal{C}$ by $\degr(u) + \degr(v)$ ascending
\For{each $\{u,v\} \in \mathcal{C}$}
    \If{$\overline{\mathbf{p}(u)\mathbf{p}(v)}$ crosses no edge of $E$}
        \State $E \gets E \cup \{\{u,v\}\}$
    \EndIf
\EndFor
\State \Return $G' = (V, E)$
\end{algorithmic}
\end{algorithm}

Sorting by minimum degree first connects the most isolated vertices before the
densely connected ones, producing a more balanced graph.

% ══════════════════════════════════════════════════════════════════════════════
\section{Tutte Embedding}
\label{sec:tutte}
% ══════════════════════════════════════════════════════════════════════════════

\subsection{Tutte's Theorem}

\begin{theorem}[Tutte, 1963~\cite{tutte1963}]
Let $G$ be a 3-connected planar graph with outer face $F_0 = (v_1, \ldots, v_k)$
fixed on a convex polygon in $\mathbb{R}^2$. If every interior vertex satisfies
the \textbf{barycentric condition}:
\begin{equation}
    \mathbf{p}(v) = \frac{1}{\degr(v)} \sum_{u \sim v} \mathbf{p}(u),
    \label{eq:barycentric}
\end{equation}
then the resulting embedding is a crossing-free straight-line drawing.
\end{theorem}

\subsection{The Linear System}

Rearranging~\eqref{eq:barycentric} for each interior vertex $v_i$ and separating
interior neighbors from boundary neighbors:
\begin{equation}
    \degr(v_i)\, \mathbf{p}(v_i)
    - \sum_{\substack{u \sim v_i \\ u \in \text{interior}}} \mathbf{p}(u)
    = \sum_{\substack{u \sim v_i \\ u \in F_0}} \mathbf{p}(u).
    \label{eq:system}
\end{equation}

In matrix form, for the $x$ and $y$ coordinates separately:
\begin{equation}
    L_{\text{int}}\, \mathbf{x} = \mathbf{b}_x, \qquad
    L_{\text{int}}\, \mathbf{y} = \mathbf{b}_y,
    \label{eq:Lx}
\end{equation}

where $L_{\text{int}}$ is the \textbf{interior Laplacian submatrix}:
\begin{equation}
    (L_{\text{int}})_{ij} =
    \begin{cases}
        \degr(v_i) & i = j, \\
        -1         & i \neq j \text{ and } v_i \sim v_j, \\
        0          & \text{otherwise.}
    \end{cases}
    \label{eq:laplacian}
\end{equation}

The right-hand side vector $\mathbf{b}_x$ has at position $i$ the sum of the
$x$-coordinates of the boundary neighbors of $v_i$.

\subsection{Properties of the Laplacian}

\begin{theorem}
The matrix $L_{\text{int}}$ defined in~\eqref{eq:laplacian} is symmetric
positive definite (SPD).
\end{theorem}

\begin{proof}
$L_{\text{int}}$ is the submatrix of the graph Laplacian induced by the interior
nodes. The full graph Laplacian of a connected graph has a single zero eigenvalue
with eigenvector proportional to the all-ones vector. Fixing the boundary nodes
(removing the corresponding rows and columns) eliminates this zero eigenvalue,
making the submatrix invertible. Positive definiteness follows from spectral
graph theory~\cite{chung1997}: all eigenvalues of $L_{\text{int}}$ are strictly
positive.
\end{proof}

Since $L_{\text{int}}$ is SPD, system~\eqref{eq:Lx} has a unique solution, computed
by \texttt{numpy.linalg.solve} via Gaussian elimination with partial pivoting
in $O(m^3)$, where $m$ is the number of interior nodes.

\subsection{Application to the Errera Graph}

The outer face $\{0, 1, 2\}$ is fixed on an equilateral triangle of radius 8:
\begin{align}
    \mathbf{p}(0) &= (0,\; 8), \notag \\
    \mathbf{p}(1) &= (-6.928,\; -4), \notag \\
    \mathbf{p}(2) &= (6.928,\; -4). \notag
\end{align}

The 14 interior nodes $\{3, 4, \ldots, 16\}$ generate a $14 \times 14$ system
with condition number $\kappa(L_{\text{int}}) \approx 18.25$, indicating a
well-conditioned system. The upper-left $4 \times 4$ block of $L_{\text{int}}$
is:

\begin{equation}
\begin{pmatrix}
 5 & -1 &  0 &  0 \\
-1 &  5 & -1 &  0 \\
 0 & -1 &  5 & -1 \\
 0 &  0 & -1 &  6
\end{pmatrix}
\end{equation}

Selected positions after solving the system:
\begin{align}
    \mathbf{p}(6)  &= (-1.7108,\; -1.0880), \notag \\
    \mathbf{p}(11) &= (-0.5338,\; -0.6087), \notag \\
    \mathbf{p}(16) &= (0.0,\; -0.7563). \notag
\end{align}

The outer face is selected by maximizing the sum of eccentricities. The
\textbf{eccentricity} of a vertex is:
\begin{equation}
    \ecc(v) = \max_{u \in V} d(v, u),
\end{equation}
where $d(v, u)$ is the shortest-path distance, computed via BFS in
$O(|V| + |E|)$. In the Errera graph: $\ecc(0) = \ecc(16) = 4$,
$\ecc(v) = 3$ for all others, giving
$\sum_{v \in \{0,1,2\}} \ecc(v) = 10$ as the maximum over all faces.

% ══════════════════════════════════════════════════════════════════════════════
\section{Topological Peeling}
\label{sec:peeling}
% ══════════════════════════════════════════════════════════════════════════════

\subsection{Face Traversal: \texorpdfstring{$\sigma \circ \alpha^{-1}$}{sigma composed with alpha inverse}}

Given the combinatorial embedding, faces are computed via dart traversal. A
\textbf{dart} is a directed edge $(v \to w)$. Two operations are defined:
\begin{align}
    \alpha(v \to w) &= (w \to v) \qquad \text{(dart reversal)}, \\
    \sigma_w(v \to w) &= \text{next dart clockwise around } w.
\end{align}

The composition $\sigma \circ \alpha^{-1}$ traverses the face adjacent to each
dart: starting from $(v_0 \to v_1)$, the sequence
\[
(v_0 \to v_1),\; (\sigma \circ \alpha^{-1})(v_0 \to v_1),\; \ldots
\]
eventually closes at $(v_0 \to v_1)$, and the visited vertices form the
face~\cite{mohar2001}. This is implemented by \texttt{emb.traverse\_face(v, w)}
in \texttt{NetworkX}.

\subsection{Peeling Algorithm}

\begin{algorithm}[H]
\caption{Topological peeling}
\label{alg:peeling}
\begin{algorithmic}[1]
\Require Maximal planar graph $G$
\Ensure Elimination order $\pi = (v_1, v_2, \ldots, v_n)$
\State $G' \gets G$, $\pi \gets ()$
\While{$V(G') \neq \emptyset$}
    \State Compute planar embedding of $G'$ (Boyer-Myrvold)
    \State $\mathcal{F} \gets$ all faces via dart traversal
    \State $F^* \gets \arg\max_{F \in \mathcal{F}} \sum_{v \in F} \ecc_{G'}(v)$
    \State $v^* \gets \arg\min_{v \in F^*} \degr_{G'}(v)$
    \State $\pi \gets \pi \cdot (v^*)$
    \State $G' \gets G' \setminus \{v^*\}$
\EndWhile
\State \Return $\pi$
\end{algorithmic}
\end{algorithm}

\subsection{Results on the Errera Graph}

The complete peeling order and degrees at the moment of elimination are:

\begin{center}
\begin{tabular}{cccc|cccc}
\toprule
\textbf{Step} & \textbf{Node} & \textbf{Outer face} & \textbf{Deg.} &
\textbf{Step} & \textbf{Node} & \textbf{Outer face} & \textbf{Deg.} \\
\midrule
 1 &  0 & $\{0,1,2\}$           & 5 & 10 &  9 & $\{9,10,11,12,13,14\}$  & 3 \\
 2 &  1 & $\{1,2,3,4,5\}$       & 4 & 11 & 10 & $\{10,11,12,13,14,15\}$ & 2 \\
 3 &  2 & $\{2,3,4,5,6,7\}$     & 3 & 12 & 11 & $\{11,12,13,14,15\}$    & 3 \\
 4 &  3 & $\{3,4,5,6,7,8\}$     & 3 & 13 & 12 & $\{12,13,14,15,16\}$    & 2 \\
 5 &  4 & $\{4,5,6,7,8,9\}$     & 3 & 14 & 13 & $\{13,14,15,16\}$       & 2 \\
 6 &  5 & $\{5,6,7,8,9,10\}$    & 2 & 15 & 14 & $\{14,15,16\}$          & 2 \\
 7 &  6 & $\{6,7,8,9,10\}$      & 4 & 16 & 15 & $\{15,16\}$             & 1 \\
 8 &  7 & $\{7,8,9,10,11,12\}$  & 3 & 17 & 16 & $\{16\}$                & 0 \\
 9 &  8 & $\{8,9,10,11,12,13\}$ & 3 &    &    &                          &   \\
\bottomrule
\end{tabular}
\end{center}

The maximum degree at elimination is 5 (step 1), consistent with
Theorem~\ref{thm:deg5}. The \textbf{degeneracy} of the graph~\cite{matula1983}
is $d(G) = \max_i \degr_{G_i}(v_i) = 5$.

% ══════════════════════════════════════════════════════════════════════════════
\section{Inverse Greedy Coloring}
\label{sec:greedy}
% ══════════════════════════════════════════════════════════════════════════════

\subsection{The Algorithm}

Given the peeling order $\pi = (v_1, \ldots, v_n)$, coloring proceeds in
reverse order $v_n, v_{n-1}, \ldots, v_1$:

\begin{algorithm}[H]
\caption{Inverse greedy coloring}
\label{alg:greedy}
\begin{algorithmic}[1]
\Require Graph $G$, peeling order $\pi = (v_1, \ldots, v_n)$
\Ensure Coloring $c: V \to \N$
\State $c \gets \{\}$
\For{$i = n, n-1, \ldots, 1$}
    \State $\mathcal{P}(v_i) \gets \{c(u) : u \in N(v_i),\; u \in \dom(c)\}$
    \State $c(v_i) \gets \min(\N \setminus \mathcal{P}(v_i))$
\EndFor
\State \Return $c$
\end{algorithmic}
\end{algorithm}

\subsection{Six-Color Guarantee}

\begin{theorem}[Six Color Theorem~\cite{heawood1890}]
Every planar graph is 6-colorable.
\end{theorem}

\begin{proof}
By induction on $|V|$. By Theorem~\ref{thm:deg5}, there exists $v$ with
$\degr(v) \leq 5$. By the inductive hypothesis, $G \setminus \{v\}$ is
6-colorable. Reinserting $v$, it has at most 5 neighbors using at most 5
distinct colors. Of the 6 available colors, at least one is free. $\square$

The inverse greedy coloring implements this argument constructively: the reverse
peeling order guarantees that when $v_i$ is colored, its already-colored
neighbors are exactly those it had when removed during peeling, and at that
moment $\degr(v_i) \leq 5$.
\end{proof}

\subsection{Results on the Errera Graph}

The inverse greedy coloring produces a valid 4-coloring of the Errera graph
directly. The complete assignment is:

\begin{center}
\begin{tabular}{ccc|ccc|ccc}
\toprule
\textbf{Node} & \textbf{Forbidden} & \textbf{Color} &
\textbf{Node} & \textbf{Forbidden} & \textbf{Color} &
\textbf{Node} & \textbf{Forbidden} & \textbf{Color} \\
\midrule
16 & $\emptyset$      & 0 &  9 & $\{0,1,2\}$ & 3 &  2 & $\{0,2,3\}$ & 1 \\
15 & $\{0\}$          & 1 &  8 & $\{1,2,3\}$ & 0 &  1 & $\{1,2,3\}$ & 0 \\
14 & $\{0,1\}$        & 2 &  7 & $\{0,1,2\}$ & 3 &  0 & $\{0,1,2\}$ & 3 \\
13 & $\{0,2\}$        & 1 &  6 & $\{0,2,3\}$ & 1 &    &             &   \\
12 & $\{0,1\}$        & 2 &  5 & $\{0,1\}$   & 2 &    &             &   \\
11 & $\{0,1,2\}$      & 3 &  4 & $\{0,2,3\}$ & 1 &    &             &   \\
10 & $\{1,3\}$        & 0 &  3 & $\{0,1,3\}$ & 2 &    &             &   \\
\bottomrule
\end{tabular}
\end{center}

The resulting coloring uses $|\{c(v) : v \in V\}| = 4$ colors and is verified
valid: $\forall \{u,v\} \in E: c(u) \neq c(v)$.

% ══════════════════════════════════════════════════════════════════════════════
\section{Kempe Chains and the Errera Counterexample}
\label{sec:kempe}
% ══════════════════════════════════════════════════════════════════════════════

\subsection{Kempe Chains}

\begin{definition}[Kempe chain~\cite{kempe1879}]
Given a partial coloring $c$ and two colors $c_i, c_j$, the
$(c_i, c_j)$-Kempe chain from $v$ is the connected component of $v$ in the
subgraph induced by vertices colored $c_i$ or $c_j$:
\begin{equation}
    K_{ij}(v) = \text{connected component of } v \text{ in }
    G[\{u : c(u) \in \{c_i, c_j\}\}].
\end{equation}
\end{definition}

A \textbf{Kempe swap} exchanges $c_i \leftrightarrow c_j$ on all nodes of
$K_{ij}(v)$. The resulting coloring remains valid because $K_{ij}(v)$ is a
connected component: no edges cross from inside $K_{ij}(v)$ to external nodes
with color $c_i$ or $c_j$~\cite{kempe1879}.

\subsection{Kempe's Error and the Errera Graph}

Kempe assumed that whenever the current outer face uses 4 colors and the next
vertex to insert has all 4 forbidden, it is always possible to reduce the outer
face to 3 colors via a Kempe swap. This assumption is false, as Heawood showed
in 1890~\cite{heawood1890}.

The Errera graph was constructed to exhibit this impossibility
concretely~\cite{errera1921}. The system analyzes all candidate chains $(c_i,
c_j)$ at each conflicting step, classifying each as:

\begin{itemize}[noitemsep]
    \item \textbf{SUCCESS}: the swap reduces the outer face to $\leq 3$ colors.
    \item \textbf{TANGLED}: the chain connects two outer-face nodes of colors
        $c_i$ and $c_j$, making the swap useless.
    \item \textbf{NO\_HELP}: valid swap but the outer face still uses 4 colors.
\end{itemize}

\begin{definition}[Tangled Kempe chain]
A $(c_i, c_j)$-chain $K_{ij}(v)$ from $v$ (with $c(v) = c_i$) is \emph{tangled}
if it contains a node $w$ on the outer face with $c(w) = c_j$. The swap
exchanges $c_i \to c_j$ at $v$ and $c_j \to c_i$ at $w$, leaving both colors
present on the outer face.
\end{definition}

At step 14 of Act 2 (insertion of node 3) in the Errera graph, all candidate
chains on the outer face are tangled. The system checks the $\binom{4}{2} = 6$
color pairs and all possible chains from outer-face nodes, finding that none
reduces the outer face from 4 to 3 colors. This is the concrete instance of the
phenomenon documented by Errera~\cite{errera1921}.

The Kittell graph~\cite{kittell1935}, with 23 vertices and $\chi(G) = 4$,
exhibits the same phenomenon at multiple steps during the Act 2 coloring.

% ══════════════════════════════════════════════════════════════════════════════
\section{CSP Backtracking with Four-Color Guarantee}
\label{sec:backtracking}
% ══════════════════════════════════════════════════════════════════════════════

\subsection{CSP Formulation}

When the greedy coloring with Kempe swaps produces more than 4 colors, the
system activates backtracking over the CSP formulation:

\begin{definition}[Constraint Satisfaction Problem, CSP]
The 4-coloring problem for $G = (V, E)$ is formulated as a CSP with:
\begin{itemize}[noitemsep]
    \item Variables: $X = \{x_v : v \in V\}$
    \item Domains: $D(x_v) = \{0, 1, 2, 3\}$ for all $v$
    \item Constraints: $x_u \neq x_v$ for all $\{u,v\} \in E$
\end{itemize}
\end{definition}

The search space is $|D|^{|V|} = 4^{17} \approx 1.7 \times 10^{10}$ for the
Errera graph. Pruning heuristics reduce the number of recursive calls to 22 in
practice.

\subsection{MRV Heuristic}

\begin{definition}[Minimum Remaining Values, MRV~\cite{russell2020}]
The MRV heuristic selects at each step the unassigned variable with the smallest
domain:
\begin{equation}
    v^* = \argmin_{v \notin \dom(c)} |\dom(x_v)|.
\end{equation}
Ties are broken by choosing the vertex of highest degree (\textit{degree
heuristic}).
\end{definition}

MRV implements the ``fail-first'' principle: by choosing the most constrained
variable, failures are detected early and more of the search tree is
pruned~\cite{haralick1980}.

\subsection{Forward Checking}

\begin{definition}[Forward Checking~\cite{haralick1980}]
Upon assigning color $c$ to vertex $v$, forward checking removes $c$ from the
domain of each uncolored neighbor $u \in N(v) \setminus \dom(c)$:
\begin{equation}
    \dom(x_u) \gets \dom(x_u) \setminus \{c\} \quad \forall u \in N(v),\; u \notin \dom(c).
\end{equation}
If any $\dom(x_u)$ becomes empty, the current assignment is inconsistent and the
algorithm backtracks immediately.
\end{definition}

\subsection{Complete Algorithm}

\begin{algorithm}[H]
\caption{Backtracking with MRV and Forward Checking}
\label{alg:backtrack}
\begin{algorithmic}[1]
\Require Graph $G = (V, E)$
\Ensure Coloring $c: V \to \{0,1,2,3\}$ or FAILURE
\State $c \gets \{\}$; $\dom(x_v) \gets \{0,1,2,3\}$ for all $v$
\Function{Backtrack}{}
    \State $v^* \gets \argmin_{v \notin \dom(c)} |\dom(x_v)|$ \hfill (MRV)
    \If{$v^* = \emptyset$} \Return \textbf{DONE} \EndIf
    \For{$\ell \in \dom(x_{v^*})$ in ascending order}
        \If{$\ell \notin \{c(u) : u \in N(v^*) \cap \dom(c)\}$}
            \State $\text{pruned} \gets \emptyset$
            \For{$u \in N(v^*) \setminus \dom(c)$} \hfill (Forward Check)
                \If{$\ell \in \dom(x_u)$}
                    \State $\dom(x_u) \gets \dom(x_u) \setminus \{\ell\}$
                    \State $\text{pruned} \gets \text{pruned} \cup \{u\}$
                    \If{$\dom(x_u) = \emptyset$}
                        \State Restore pruned; \textbf{continue} to next $\ell$
                    \EndIf
                \EndIf
            \EndFor
            \State $c(v^*) \gets \ell$
            \If{\Call{Backtrack}{} = \textbf{DONE}} \Return \textbf{DONE} \EndIf
            \State $c(v^*) \gets \text{undefined}$; restore domains of pruned
        \EndIf
    \EndFor
    \State \Return \textbf{FAILURE}
\EndFunction
\end{algorithmic}
\end{algorithm}

\subsection{Search Trace Visualization}

The system records each event of the search tree for frame-by-frame visualization:

\begin{center}
\begin{tabular}{llp{7cm}}
\toprule
\textbf{Event} & \textbf{Frame} & \textbf{Meaning} \\
\midrule
\texttt{ASSIGN}   & Color fixed, neighbors pruned
    & Successful assignment; node shown with thick white border \\
\texttt{FAIL}     & Node in dark red
    & Domain wipeout in some neighbor; must backtrack \\
\texttt{UNASSIGN} & Node decolored
    & Undoes assignment and restores domains \\
\texttt{DONE}     & All nodes colored
    & Solution found \\
\bottomrule
\end{tabular}
\end{center}

\subsection{Results on the Errera Graph}

\begin{center}
\begin{tabular}{ll}
\toprule
\textbf{Parameter} & \textbf{Value} \\
\midrule
Recursive calls       & 22 \\
\texttt{ASSIGN} events & 21 \\
\texttt{FAIL} / \texttt{UNASSIGN} events & 4 each \\
Colors used           & 4 \\
Valid coloring         & Yes \\
Full search space     & $4^{17} \approx 1.7 \times 10^{10}$ \\
\bottomrule
\end{tabular}
\end{center}

The coloring found by backtracking is:
$\{6{:}0,\, 7{:}1,\, 1{:}2,\, 2{:}0,\, 8{:}2,\, 0{:}1,\, 3{:}3,\, 5{:}3,\,
9{:}0,\, 4{:}2,\, 10{:}1,\, 11{:}3,\, 12{:}2,\, 15{:}2,\, 13{:}0,\, 16{:}1,\,
14{:}3\}$, differing from the greedy coloring in 15 of 17 nodes.

% ══════════════════════════════════════════════════════════════════════════════
\section{The Four Color Theorem: Historical Context}
\label{sec:history}
% ══════════════════════════════════════════════════════════════════════════════

\subsection{Timeline}

\begin{description}[leftmargin=2cm]
    \item[1852] Francis Guthrie formulates the conjecture while coloring
        maps~\cite{wilson2002}.
    \item[1879] Alfred Kempe proposes a proof based on chain
        interchange~\cite{kempe1879}.
    \item[1890] Percy Heawood identifies the flaw in Kempe's proof and proves
        the Five Color Theorem~\cite{heawood1890}.
    \item[1921] Alfred Errera constructs a concrete planar graph exhibiting the
        failure point of Kempe's algorithm~\cite{errera1921}.
    \item[1935] I. Kittell constructs a second counterexample to Kempe's
        algorithm~\cite{kittell1935}.
    \item[1976] Appel and Haken prove the Four Color Theorem by computationally
        verifying 1936 reducible configurations~\cite{appelhaken1977}.
    \item[1997] Robertson, Sanders, Seymour, and Thomas simplify the proof to
        633 configurations with an independent verifier~\cite{robertson1997}.
    \item[2005] Georges Gonthier produces the first mechanically verified formal
        proof in the Coq proof assistant~\cite{gonthier2008}.
\end{description}

\subsection{Gonthier's Formal Verification}

Gonthier's contribution~\cite{gonthier2008} shifts the object of trust: rather
than trusting a 500,000-line verification program, one trusts only the Coq
kernel (approximately 3,000 lines), whose correctness can be established
independently.

The formalism uses \textbf{combinatorial hypermaps}: an algebraic representation
of a planar embedding based on two permutations over a set of darts. If $H$
denotes the dart set and $\sigma, \alpha: H \to H$ are permutations with
$\alpha^2 = \mathrm{id}$, then $(H, \sigma, \alpha)$ is a hypermap representing
the embedding. Faces are orbits of $\sigma \circ \alpha$, vertices are orbits of
$\sigma$, and edges are orbits of $\alpha$~\cite{gonthier2008}.

This is isomorphic to the combinatorial embedding used by \texttt{NetworkX}
(implemented as a dict-of-dicts storing the cyclic neighbor order), though
Gonthier's presentation is more formal and amenable to mechanical verification.

\subsection{What the Code Implements and What It Does Not}

The implemented system constructively demonstrates that the Errera and Kittell
graphs are 4-colorable. It does not prove the Four Color Theorem for arbitrary
planar graphs, which requires:

\begin{enumerate}[noitemsep]
    \item The \textbf{discharging method}: assign charge $\mu(v) = 6 - \degr(v)$
        to each vertex (the total sum is 12 by Euler) and redistribute it via
        discharge rules that yield a contradiction if the graph contains no
        reducible configuration.
    \item \textbf{Reducibility}: for each of the 633 configurations, show that
        any coloring of the outer ring can be extended to the interior (possibly
        with Kempe swaps).
    \item \textbf{Unavoidability}: show that the configuration set is unavoidable,
        i.e., every planar graph contains at least one~\cite{robertson1997}.
\end{enumerate}

% ══════════════════════════════════════════════════════════════════════════════
\section{Computational Complexity}
\label{sec:complexity}
% ══════════════════════════════════════════════════════════════════════════════

Table~\ref{tab:complexity} summarizes the complexity of each pipeline stage.

\begin{table}[H]
\centering
\caption{Complexity of each pipeline stage}
\label{tab:complexity}
\begin{tabular}{llll}
\toprule
\textbf{Stage} & \textbf{Algorithm} & \textbf{Complexity} & \textbf{Reference} \\
\midrule
Graph loading        & Adjacency list insertion  & $O(|E|)$             & --- \\
Planarity check      & Boyer-Myrvold             & $O(|V|)$             & \cite{boyermyrvold1999} \\
Maximalization       & Brute force + Shapely     & $O(|V|^2 |E|)$       & --- \\
Tutte embedding      & Gaussian elim.\ ($m \times m$) & $O(|V|^3)$      & \cite{tutte1963} \\
BFS (eccentricity)   & BFS from each vertex      & $O(|V|(|V|+|E|))$    & --- \\
Face traversal       & Dart traversal            & $O(|F|)$             & \cite{mohar2001} \\
Full peeling         & Peeling $\times$ Boyer-Myrvold & $O(|V|^2)$      & --- \\
Greedy coloring      & Greedy over adjacency     & $O(|V|+|E|)$         & --- \\
CSP backtracking (worst) & MRV + FC              & $O(4^{|V|})$         & \cite{haralick1980} \\
CSP backtracking (practice) & Planar structure  & $O(|V|)$             & --- \\
\bottomrule
\end{tabular}
\end{table}

The theoretical bottleneck of backtracking is exponential in the worst case, but
in practice the structure of planar graphs causes MRV + FC to solve the problem
in a few dozen recursive calls. The Errera graph required 22 calls instead of
$4^{17}$.

% ══════════════════════════════════════════════════════════════════════════════
\section{Implementation}
\label{sec:implementation}
% ══════════════════════════════════════════════════════════════════════════════

\subsection{Dependencies}

\begin{center}
\begin{tabular}{lll}
\toprule
\textbf{Library} & \textbf{Version} & \textbf{Primary use} \\
\midrule
\texttt{networkx}  & $\geq$ 3.0  & Graph structure, planar embedding, BFS \\
\texttt{numpy}     & $\geq$ 1.24 & Tutte linear system, linear algebra \\
\texttt{scipy}     & $\geq$ 1.10 & \texttt{ConvexHull} (legacy visual aid) \\
\texttt{shapely}   & $\geq$ 2.0  & Segment intersection test \\
\texttt{matplotlib}& $\geq$ 3.7  & Visualization and GIF animation \\
\texttt{pillow}    & $\geq$ 9.0  & GIF export \\
\bottomrule
\end{tabular}
\end{center}

\subsection{Code Structure}

The main file \texttt{onion\_peeling.py} organizes the pipeline into eight
numbered sections, each implementing one mathematical stage:

\begin{enumerate}[noitemsep]
    \item \textbf{Loading}: JSON reading and graph construction.
    \item \textbf{Maximalization}: Algorithm~\ref{alg:maximal}.
    \item \textbf{Tutte embedding}: Laplacian system, Section~\ref{sec:tutte}.
    \item \textbf{Outer face}: face traversal and eccentricity.
    \item \textbf{Peeling}: Algorithm~\ref{alg:peeling}.
    \item \textbf{Greedy coloring}: Algorithm~\ref{alg:greedy}.
    \item \textbf{Kempe chains}: chain analysis, Section~\ref{sec:kempe}.
    \item \textbf{Backtracking}: Algorithm~\ref{alg:backtrack}.
    \item \textbf{Animation}: three-act GIF generation.
\end{enumerate}

The system produces two main outputs:
\begin{itemize}[noitemsep]
    \item \texttt{descomposicion\_grafo.gif}: animated process on the input graph.
    \item \texttt{explainer.gif}: 40-frame pedagogical GIF documenting each
        mathematical stage.
\end{itemize}

% ══════════════════════════════════════════════════════════════════════════════
\section{Layer-by-Layer Tab\"{u} Coloring}
\label{sec:tabu}
% ══════════════════════════════════════════════════════════════════════════════

\subsection{Motivation}

The greedy inverse coloring of Section~\ref{sec:greedy} inserts vertices one at
a time in reverse peel order. When it produces more than four colors, the system
falls back to CSP backtracking (Section~\ref{sec:backtracking}), which is
guaranteed to find a four-coloring but may explore a large portion of the search
space. The question motivating this section is: \emph{can the coloring be
organized so that backtracking is never needed?}

The key observation is that the topological peeling naturally groups vertices
into \textbf{layers}: all vertices on the outer face at a given moment form one
layer, and removing that layer exposes the next. Coloring an entire layer at
once, before moving inward, provides more global information at each decision
point than the node-by-node greedy approach.

\subsection{Algorithm}

\begin{definition}[Peeling layers]
The peeling layers $L_0, L_1, \ldots, L_k$ of a planar graph $G$ are defined
iteratively: $L_0$ is the outer face of $G$; for $i > 0$, $L_i$ is the outer
face of $G \setminus (L_0 \cup \cdots \cup L_{i-1})$.
\end{definition}

For the Errera graph the layers are:
\[
    L_0 = \{0,1,2\}, \quad
    L_1 = \{3,4,5,6,7,8\}, \quad
    L_2 = \{9,10,11,12,13,14\}, \quad
    L_3 = \{15,16\}.
\]

\begin{definition}[Tab\"{u} constraint]
When coloring a vertex $v \in L_i$, the \textbf{tab\"{u}} set of $v$ is the set
of colors already assigned to its direct neighbors in previously colored layers:
\begin{equation}
    \tau(v) = \{c(u) : u \in N(v),\; u \in L_0 \cup \cdots \cup L_{i-1}\}.
    \label{eq:tabu}
\end{equation}
\end{definition}

The coloring rule is identical to the greedy rule of Section~\ref{sec:greedy}
— assign the minimum color not in $\tau(v)$ — but the \emph{scope} differs:
here $\tau(v)$ is computed from all previously colored layers, whereas in the
node-by-node greedy it includes only the previously inserted nodes in peel order.
Within each layer, nodes are colored in ascending ID order.

\begin{algorithm}[H]
\caption{Layer-by-layer tab\"{u} coloring}
\label{alg:tabu}
\begin{algorithmic}[1]
\Require Planar graph $G$
\Ensure Coloring $c: V \to \N$
\State $c \gets \{\}$,\quad $i \gets 0$
\While{$G \neq \emptyset$}
    \State $L_i \gets \text{outer face of } G$ \Comment{face traversal, max ecc sum}
    \For{$v \in L_i$ in ascending order}
        \State $\tau(v) \gets \{c(u) : u \in N(v),\; u \in \dom(c)\}$
        \State $c(v) \gets \min(\N \setminus \tau(v))$
    \EndFor
    \State $G \gets G \setminus L_i$,\quad $i \gets i + 1$
\EndWhile
\State \Return $c$
\end{algorithmic}
\end{algorithm}

\subsection{Relationship to Greedy and to Backtracking}

The tab\"{u} algorithm is strictly more informed than the node-by-node greedy:
when a vertex $v \in L_i$ is colored, \emph{all} of its neighbors in
$L_0, \ldots, L_{i-1}$ are already colored, whereas in the greedy they may not
be (depending on the peel order within earlier layers). This means $\tau(v)$
is at least as large as the forbidden set in the node-by-node greedy, so the
tab\"{u} algorithm makes more constrained — and therefore more globally
consistent — decisions.

The tab\"{u} constraint~\eqref{eq:tabu} is equivalent to the \emph{initial}
forward-checking propagation in the CSP backtracking: before assigning any color
to $v$, backtracking would also eliminate the colors of all colored neighbors
from $v$'s domain. The difference is that the tab\"{u} algorithm never revises
a decision, whereas backtracking undoes assignments when a domain becomes empty.

\subsection{Results on the Errera Graph}

Applying Algorithm~\ref{alg:tabu} to the Errera graph:

\begin{center}
\begin{tabular}{clll}
\toprule
\textbf{Layer} & \textbf{Nodes} & \textbf{Tab\"{u} sets} & \textbf{Colors assigned} \\
\midrule
$L_0$ & $\{0,1,2\}$
    & $\emptyset,\;\{0\},\;\{0,1\}$
    & $0,\;1,\;2$ \\
$L_1$ & $\{3,4,5,6,7,8\}$
    & $\{0,2\},\;\{0,1\},\;\{0,1,2\},\;\{1,3\},\;\{0,1,2\},\;\{1,2,3\}$
    & $1,\;2,\;3,\;0,\;3,\;0$ \\
$L_2$ & $\{9,10,11,12,13,14\}$
    & $\{0,1,2\},\;\{0,2,3\},\;\{0,1\},\;\{0,2,3\},\;\{0,1,3\},\;\{0,2,3\}$
    & $3,\;1,\;2,\;1,\;2,\;1$ \\
$L_3$ & $\{15,16\}$
    & $\{1,2,3\},\;\{0,1,2\}$
    & $0,\;3$ \\
\bottomrule
\end{tabular}
\end{center}

The algorithm produces a valid 4-coloring \textbf{without backtracking}. The
maximum tab\"{u} set size is 3 (occurring at nodes 5, 7, 8 in $L_1$, and
several nodes in $L_2$), consistent with the bound $|\tau(v)| \leq 5$ that
follows from the density argument of Theorem~\ref{thm:deg5}.

\subsection{Comparison with Other Methods}

\begin{center}
\begin{tabular}{lcccc}
\toprule
\textbf{Method} & \textbf{Colors} & \textbf{Backtracks} & \textbf{Assignments} & \textbf{Complexity} \\
\midrule
Greedy (node-by-node)      & 4 & 0  & 17 & $O(V+E)$ \\
Greedy + Kempe swaps       & 5 & 0  & 17 & $O(V(V+E))$ \\
CSP Backtracking (MRV+FC)  & 4 & 4  & 21 & $O(4^V)$ worst \\
Tab\"{u} layer-by-layer    & 4 & 0  & 17 & $O(V+E)$ \\
\bottomrule
\end{tabular}
\end{center}

The tab\"{u} algorithm matches the result of CSP backtracking (4 colors, valid)
without any backtracking, and runs in $O(V+E)$ — the same asymptotic complexity
as the greedy. The key open question is whether this holds for all planar graphs
or only for structured cases like Errera.

% ══════════════════════════════════════════════════════════════════════════════
\section{Bi-Layer Coloring: Failure Detection and Closure}
\label{sec:bilayer}
% ══════════════════════════════════════════════════════════════════════════════

\subsection{Motivation and Overview}

De Ita Luna and Marcial-Romero~\cite{deita2024} propose the Bi-Layer algorithm,
which extends the outer-face greedy approach with two key mechanisms: (1)
\emph{failure configuration detection} — recognizing local graph patterns that
would require a fifth color before they are reached — and (2) the \emph{closure
operator} $T((v_i, a))$ — simulating the transitive propagation of forced
colorings when color $a$ is tentatively assigned to vertex $v_i$, in order to
evaluate the downstream consequences of each color choice.

\subsection{Failure Configurations}

\begin{definition}[Basic failure configurations~\cite{deita2024}]
Given a partial coloring $c$ and tabu sets $\tau(v) = \{c(u) : u \sim v,\,
u \in \dom(c)\}$:
\begin{itemize}[noitemsep]
    \item \textbf{Failed vertex} (Fig.~1a): $|\tau(v)| = 4$ — all four colors forbidden.
    \item \textbf{Failed edge} (Fig.~1b): $\{v_1, v_2\} \in E$,
        $\tau(v_1) = \tau(v_2)$, $|\tau(v_1)| = 3$.
    \item \textbf{Failed face} (Fig.~1c): triangle $v_1$-$v_2$-$v_3$,
        $\tau(v_1) = \tau(v_2) = \tau(v_3)$, $|\tau(v_1)| = 2$.
\end{itemize}
\end{definition}

Any of these configurations guarantees that the coloring cannot be completed
with four colors. A \textbf{failed wheel} $W_4$ consists of two potential
failed edges with complementary tabu sets sharing a common hub.

\subsection{Potential Failure Configurations}

\begin{definition}[Potential failed edge (pfe)~\cite{deita2024}]
An edge $\{v_1, v_2\}$ is a potential failed edge if:
\[
    \tau(v_1) = \tau(v_2),\quad |\tau(v_1)| = 2,\quad
    \exists w \in G_i : v_1, v_2 \in N(w).
\]
\end{definition}

The algorithmic strategy for pfe: color the endpoints $v_1, v_2$
\emph{before} their common neighbor $w$. This prevents the failed edge from
materializing when $w$ is subsequently colored.

\begin{definition}[Potential failed face (pff)~\cite{deita2024}]
A face $v_1$-$v_2$-$v_3$ is a potential failed face if:
\[
    \left(\bigcap_{v \in \{v_1,v_2,v_3\}} \tau(v)\right) \neq \emptyset,\quad
    |\tau(v_1)| = 2,\quad
    \exists w \in G_i : v_2, v_3 \in N(w),\; v_1 \notin N(w).
\]
\end{definition}

Strategy for pff: assign to $w$ a color not appearing in the tabu of any
vertex at distance 2 from $w$ (i.e., not in $\tau(N_2(w))$).

\subsection{The Closure Operator}

\begin{definition}[Closure $T((v_i, a))$~\cite{deita2024}]
Given partial coloring $c$ and a tentative assignment $c(v_i) = a$, the
closure propagates forced colorings transitively:
\begin{enumerate}[noitemsep]
    \item Initialize $\text{forced} = \{v_i \mapsto a\}$.
    \item For each $u \notin \dom(c)$ not yet in \emph{forced}: compute
        $\tau'(u) = \tau(u) \cup \{forced[nb] : nb \in N(u) \cap \text{forced}\}$.
    \item If $|\text{FOUR} \setminus \tau'(u)| = 0$ $\Rightarrow$ INCONSISTENT.
    \item If $|\text{FOUR} \setminus \tau'(u)| = 1$ $\Rightarrow$ add forced assignment
        and continue propagation.
\end{enumerate}
If the closure produces an inconsistency (failed vertex, edge, or face),
color $a$ is rejected for $v_i$.
\end{definition}

\subsection{Vertex and Color Selection}

\textbf{Vertex selection} (hierarchical, from outer face):
\begin{enumerate}[noitemsep]
    \item Vertices involved in potential failed edges.
    \item Vertices with complementary tabu sets at distance 2.
    \item Vertex of maximum degree in the outer face.
\end{enumerate}

\textbf{Color selection}: for each candidate color $a \in \text{FOUR} \setminus \tau(v_i)$:
\begin{enumerate}[noitemsep]
    \item Run $T((v_i, a))$; if inconsistent, reject $a$.
    \item Build virtual graph $G_{i+1,a} = G_i \setminus \text{Proj}(T((v_i,a)))$.
    \item Score $(G_{i+1,a}, c')$ by: (a) future failed vertices, (b) failed
        edges, (c) failed faces, (d) pfe count, (e) complementary tabu pairs.
    \item Choose $a$ minimizing the score lexicographically.
\end{enumerate}

\textbf{Stack}: vertices satisfying $\deg_{G_i}(v) + |\tau(v)| \leq 3$ are
pushed to a stack and colored last (greedily); they can always be colored
correctly once all their neighbors are colored.

\subsection{Implementation Status}

The implementation in \texttt{bilayer\_coloring.py} covers: closure $T((v_i,a))$
with one-level propagation from forced neighbors; failure detection (failed
vertices, edges, faces); pfe detection; complementary-tabu detection; the stack
mechanism; hierarchical vertex selection; and closure-based color scoring.

A known limitation of the current implementation is that the closure does not
perform a global pre-pass to detect vertices already forced (with $|\text{remaining}|=1$)
before propagating from $v_i$. This can allow failure configurations to form
across non-adjacent vertices (as observed in the Errera graph at step 9, where
nodes 3, 4, and 5 simultaneously reach $|\text{remaining}|=1$ without being
direct neighbors of the currently colored vertex). Addressing this requires the
full transitive-closure implementation described in~\cite{deita2024}, including
the reflexive-transitive closure process for anticipating multi-step consequences.

\subsection{Results}

\begin{center}
\begin{tabular}{lcccc}
\toprule
\textbf{Graph} & \textbf{Tabu colors} & \textbf{Bi-Layer colors} & \textbf{Valid} & \textbf{Steps} \\
\midrule
Errera (17)       & 4 & 4 & Partial* & 12 \\
Kittell (23)      & 5 & 4 & Yes      & 16 \\
Dodecahedron (20) & 5 & 4 & Yes      & 12 \\
Icosahedron (12)  & 4 & 4 & Yes      &  6 \\
Octahedron (6)    & 3 & 3 & Yes      &  6 \\
Wheel $W_8$ (8)   & 4 & 4 & Yes      &  8 \\
Grid $4\times4$ (16) & 5 & 4 & Yes   &  6 \\
Sunflower (13)    & 4 & 4 & Yes      & 13 \\
\bottomrule
\end{tabular}
\end{center}

\small{*Errera: one edge conflict due to the global-pass limitation in the closure.
The Bi-Layer correctly identifies the pfe configurations but requires the full
transitive closure to prevent the downstream failure.}

% ══════════════════════════════════════════════════════════════════════════════
\section{Forest Decomposition Structure of the Tab\"{u} Ordering}
\label{sec:forest}
% ══════════════════════════════════════════════════════════════════════════════

\subsection{Motivation}

Computational analysis of the tab\"{u} coloring algorithm
(Section~\ref{sec:tabu}) reveals a striking structural property: the
outer-face tab\"{u} ordering implicitly partitions the vertex set into four
classes with precise combinatorial structure. This observation connects the
algorithm to classical graph arboricity theory and suggests a new approach to
proving the Four Color Theorem constructively.

\subsection{Phase Decomposition}

\begin{definition}[Backward set and phase]
Given the tab\"{u} ordering $\pi = (v_1, \ldots, v_n)$, define the
\emph{backward color set} of $v_i$ as
\[
    BC(v_i) = \bigl\{c(v_j) : j < i,\; v_j \sim v_i\bigr\},
\]
the set of distinct colors already assigned to $v_i$'s neighbors when $v_i$ is
colored. The \emph{phase} of $v_i$ is $\phi(v_i) = |BC(v_i)|$.
The \emph{phase sets} are $F_k = \{v_i : \phi(v_i) = k\}$.
\end{definition}

\subsection{Proven Structural Properties}

\begin{theorem}[Phase independence]
\label{thm:phase-indep}
Under any tabu coloring ordering, $F_k$ is an independent set for all $k \geq 2$.
\end{theorem}

\begin{proof}
Let $u, w \in F_k$ with $k \geq 2$, with $u$ colored before $w$ (so $u \to w$
is a backward edge). Since $w \in F_k$, all backward neighbors of $w$ have at
most $k$ distinct colors. But $u \in F_k$ means $c(u)$ is some color $\gamma$.
Additionally, $BC(u)$ already contains $k$ distinct colors, all present in
$w$'s backward neighborhood (as $u$'s previous backward neighbors are also
backward neighbors of $w$ along edges of the graph). Hence $|BC(w)| \geq k+1$,
placing $w$ in $F_{k+1}$, not $F_k$. Contradiction.

More precisely for $k=2$: $u \in F_2$ means $|BC(u)| = 2$, so $c(u) \in
\{0,1,2\} \setminus BC(u)$, giving $c(u) = $ the third color. If $w \sim u$
(backward), then $w$ sees $c(u)$ plus the two colors in $BC(u)$, giving
$|BC(w)| \geq 3$, so $w \in F_{\geq 3}$. By symmetry $u \notin F_2$ if $w \in
F_2$. Hence $F_2$ is independent. The argument for $k \geq 3$ is identical.
\end{proof}

\begin{corollary}
$F_2$, $F_3$, $F_4$ are each independent sets: each needs exactly one new
color. When $F_4 = \emptyset$, the algorithm uses at most 4 colors.
\end{corollary}

\subsection{The Phase-1 Forest Property}

\begin{definition}[Phase-1 backward subgraph]
The \emph{phase-1 backward subgraph} $B_1$ has vertex set $F_1$ and edge set
$\{(u,w) : u,w \in F_1, u \sim w, u \text{ colored before } w\}$.
\end{definition}

\begin{lemma}[k=3 impossible]
\label{lem:k3}
The backward subgraph $B_1$ contains no triangle.
\end{lemma}

\begin{proof}
Suppose $u, v, w \in F_1$ with $w$ colored first, then $v$, then $u$, and
$w \sim v \sim u \sim w$ (all pairs adjacent). Since $v \in F_1$: $BC(v) =
\{c(w)\}$ (exactly one color). Since $u \in F_1$: $BC(u)$ has exactly one
color. But $u \sim v$ and $u \sim w$, so $BC(u) \supseteq \{c(v), c(w)\}$.
Since $v \in F_1$ with $v \sim w$: $c(v) \neq c(w)$ (valid coloring). Thus
$|BC(u)| \geq 2$, placing $u$ in $F_{\geq 2}$. Contradiction.
\end{proof}

\begin{lemma}[Even cycles only]
\label{lem:even}
Any cycle in $B_1$ has even length.
\end{lemma}

\begin{proof}
Nodes in $B_1$ use only colors $\{0,1\}$ (they see exactly one color backward,
so they take the other). Adjacent nodes in $B_1$ have opposite colors ($u \in
F_1$ with backward neighbor $w$: $c(u) \neq c(w)$ since the coloring is
valid). Hence $B_1$ is 2-colorable with $\{0,1\}$, meaning all cycles in $B_1$
have even length.
\end{proof}

By Lemma~\ref{lem:k3}, the minimum cycle length in $B_1$ is $\geq 4$. By
Lemma~\ref{lem:even}, it must be even. So the minimum possible cycle has
length~4.

\subsection{Main Theorem and Proof}

\begin{theorem}[F1 Forest Theorem]
\label{thm:f1forest}
Let $G$ be a maximal planar graph (every face is a triangle). Under the
outer-face tab\"{u} ordering, the phase-1 backward subgraph $B_1$ is a forest.
\end{theorem}

The proof proceeds through two lemmas.

\begin{lemma}[E-before-I ordering]
\label{lem:ebeforei}
Let $C = v_1\text{-}v_2\text{-}\cdots\text{-}v_k\text{-}v_1$ be a cycle in
$B_1$. For any edge $\{v_i, v_{i+1}\}$ in $C$ and any vertex $w \notin C$
adjacent to both $v_i$ and $v_{i+1}$: if $w$ lies in the exterior region $E$
of $C$ in the planar embedding, then $w$ is colored before $v_i$.
\end{lemma}

\begin{proof}
Suppose for contradiction that $w$ is not colored before $v_i$. There are two
cases.

\textbf{Case 1}: $w$ is colored after $v_{i+1}$.
The tab\"{u} algorithm colors from the outer face inward. Nodes colored after
all cycle nodes are topologically interior to $C$ in the planar sense: they are
unreachable from the outer boundary without passing through $C$. Thus $w$ lies
in the interior region $I$ of $C$. But $w \in E$ by hypothesis. Contradiction.

\textbf{Case 2}: $w$ is colored after $v_i$ but before $v_{i+1}$.
Then $w$ is a backward neighbor of $v_{i+1}$ but not of $v_i$.
Since $v_{i+1} \in F_1$, all its backward neighbors have the same color:
\[
    c(w) = c_{\mathrm{back}}(v_{i+1}) = 1 - c(v_i) = 1-\alpha.
\]
But $w \sim v_i$ and $c(v_i) = \alpha$, so the coloring assigns $w$ a color
not in $\{\alpha\}$. Since $c(w) = 1-\alpha \neq \alpha$, no coloring conflict
arises from this. However, when $w$ is colored it is adjacent to $v_i$ (already
colored $\alpha$), so $w$'s tab\"{u} includes $\alpha$. The greedy rule assigns
$w$ the minimum color not in its tab\"{u}. If $w$'s only tab\"{u} color (from
$v_i$) is $\alpha$, then $c(w) = 1-\alpha$.

Now $v_{i+1}$'s backward neighbors include $v_i$ (color $\alpha$) and $w$
(color $1-\alpha$). These are two \emph{distinct} colors, so $|BC(v_{i+1})| \geq 2$,
placing $v_{i+1} \in F_{\geq 2}$. This contradicts $v_{i+1} \in F_1$.

In both cases we reach a contradiction. Therefore $w$ must be colored before
$v_i$.
\end{proof}

\begin{lemma}[Triangular E-face]
\label{lem:triface}
In a maximal planar graph $G$, for every edge $\{u, v\} \in E(G)$, there
exists a vertex $w$ such that $u$-$v$-$w$ is a triangular face in the planar
embedding. For any cycle $C$ containing edge $\{u,v\}$, one of the two
triangular faces of $\{u,v\}$ lies in the exterior region $E$ of $C$; call
its third vertex $w_E$.
\end{lemma}

\begin{proof}
In a maximal planar graph, every face is a triangle, so every edge borders
exactly two triangular faces. The cycle $C$ divides the plane into two open
regions; the edge $\{u,v\}$ lies on the boundary of $C$, so one face is in $I$
and one in $E$. The third vertex $w_E$ of the $E$-face exists by maximality.
\end{proof}

\begin{proof}[Proof of Theorem~\ref{thm:f1forest}]
Suppose for contradiction that $B_1$ contains a cycle
$C = v_1\text{-}v_2\text{-}\cdots\text{-}v_k\text{-}v_1$ (with vertices
colored in this backward order, i.e., $v_k$ colored first, then $v_{k-1}$,
\ldots, then $v_1$ last).

By Lemma~\ref{lem:k3}, no triangle exists in $B_1$, so $k \geq 4$. By
Lemma~\ref{lem:even}, $k$ is even and the colors alternate:
$c(v_i) \in \{0,1\}$ with $c(v_i) + c(v_{i+1}) = 1$ for all $i$ (indices
mod $k$). Write $c(v_i) = \alpha_i$ with $\alpha_i \in \{0,1\}$ and
$\alpha_{i+1} = 1 - \alpha_i$.

Take any consecutive pair $v_i$, $v_{i+1}$ in $C$ (so $v_{i+1}$ is colored
before $v_i$ in the backward order). By Lemma~\ref{lem:triface}, there exists
$w_E$ adjacent to both $v_i$ and $v_{i+1}$, lying in the exterior region $E$
of $C$.

By Lemma~\ref{lem:ebeforei}, $w_E$ is colored before $v_{i+1}$ (the earlier
of the two), hence $w_E$ is a backward neighbor of both $v_i$ and $v_{i+1}$.

Since $v_{i+1} \in F_1$: all backward neighbors of $v_{i+1}$ have the same
color, so
\[
    c(w_E) = c_{\mathrm{back}}(v_{i+1}) = 1 - \alpha_{i+1} = \alpha_i.
\]
Since $v_i \in F_1$: all backward neighbors of $v_i$ have the same color, so
\[
    c(w_E) = c_{\mathrm{back}}(v_i) = 1 - \alpha_i.
\]
But $\alpha_i \neq 1 - \alpha_i$ for $\alpha_i \in \{0,1\}$. This is a
contradiction.

Therefore $B_1$ contains no cycle, i.e., $B_1$ is a forest.
\end{proof}

\begin{corollary}[4-color bound without backtracking]
\label{cor:4color}
Let $G$ be a maximal planar graph. The outer-face tab\"{u} algorithm produces a
valid coloring of $G$ using at most $1 + 1 + 1 + 1 = 4$ colors: color $0$ for
$F_0$, colors $\{0,1\}$ for $F_1$ (a forest, hence 2-colorable), color $2$ for
$F_2$, and color $3$ for $F_3$ — provided $F_4 = \emptyset$.
\end{corollary}

\begin{remark}[On the condition $F_4 = \emptyset$]
Theorem~\ref{thm:f1forest} proves that $F_1$ is a forest, eliminating one
source of failure. Whether $F_4 = \emptyset$ always holds — i.e., whether the
tab\"{u} ordering guarantees that no vertex ever sees 4 distinct backward colors
— is equivalent to the Four Color Theorem itself. The theorem above provides a
structural characterization of \emph{why} the algorithm works as well as it
does, and reduces the 4-coloring question to a single combinatorial property
($F_4 = \emptyset$) of the ordering.
\end{remark}

\subsection{Computational Verification}

Theorem~\ref{thm:f1forest} was verified on:
\begin{itemize}[noitemsep]
    \item All named graphs: Errera, Kittell, dodecahedron, icosahedron,
        octahedron, wheel graphs, grids — $B_1$ is a forest in all cases.
    \item 474 random Delaunay triangulations: 3 cases where $B_1$ had a
        cycle were all non-maximal planar graphs, consistent with the
        theorem's hypothesis of maximality.
\end{itemize}

\textbf{Important correction.} An earlier version of this document reported that
2\,000 random triangulations with triangular outer face produced zero failures.
That test used a graph generation method that inadvertently produced planar
triangulations where the outer face happened to be a triangle, but not all faces
were triangular (non-maximal). On correct maximal planar triangulations with
triangular outer face, the plain tab\"{u} algorithm achieves 4 colors in
approximately 45\% of cases. This does not contradict
Theorem~\ref{thm:f1forest}: the theorem guarantees $F_1$ is a forest, but
$F_1$ being a forest is necessary, not sufficient, for 4-coloring. The failure
cases all correspond to $F_4 \neq \emptyset$.

\subsection{Attempts to Eliminate $F_4$}

Three strategies were tested to prevent $F_4$ from appearing:

\textbf{Interrupt strategy.} Before selecting from the outer face, check if any
uncolored node has $|\tau(v)| = 3$ (one color left). If so, color it immediately.
This prevents the node from reaching $|\tau| = 4$ due to subsequent neighbor
colorings. Result: 32\% 4-colored on random Delaunay (vs 22\% plain tabu).

\textbf{Recolor strategy.} When a node $v$ has $|\tau(v)| = 4$ (no color
available), attempt to recolor a neighbor $u$: find a new color for $u$ that
does not conflict with $u$'s other neighbors, freeing $u$'s original color for
$v$. Result: 51\% 4-colored, always valid, no backtracking.

\textbf{Kempe chain strategy.} When stuck, attempt a Kempe chain swap on
neighbor colors. Result: 33\% 4-colored.

\begin{center}
\begin{tabular}{lcc}
\toprule
\textbf{Strategy} & \textbf{4-colored} & \textbf{Always valid} \\
\midrule
Plain tab\"{u}             & 22\% & Yes \\
Tab\"{u} + interrupt       & 32\% & Yes \\
Tab\"{u} + Kempe swap      & 33\% & Yes \\
Tab\"{u} + recolor         & 51\% & Yes \\
\bottomrule
\end{tabular}
\end{center}

None of these strategies achieves 100\%. The fundamental reason: preventing
$F_4$ from appearing is equivalent to showing that no planar graph has a vertex
ordering where some node sees 4 distinct backward colors — which is exactly the
Four Color Theorem stated algorithmically. The recoloring trick is essentially
a local Kempe swap, the same operation at the heart of Kempe's original (flawed)
proof attempt.

\subsection{Connection to Arboricity}

\begin{theorem}[Nash-Williams, 1961~\cite{nashwilliams1961}]
Every planar graph satisfies $a(G) \leq 3$: its edges decompose into at most 3
forests.
\end{theorem}

The phase decomposition provides a \emph{vertex} analogue: $V = F_0 \cup F_1
\cup F_2 \cup F_3$, where $F_0 \cup F_1$ induces a forest (Theorem~\ref{thm:f1forest})
and $F_2, F_3$ are independent sets (Theorem~\ref{thm:phase-indep}). This
4-partition directly gives the 4-coloring when $F_4 = \emptyset$. The
connection to arboricity is that the tab\"{u} ordering implicitly builds a
2-forest cover of the vertex set, corresponding to Nash-Williams' edge forest
decomposition with arboricity 2 for the backward subgraph.

% ══════════════════════════════════════════════════════════════════════════════
\section{Future Work}
\label{sec:future}
% ══════════════════════════════════════════════════════════════════════════════

\subsection{Backtracking-Free Coloring --- Open Questions}

The layer-by-layer tab\"{u} algorithm of Section~\ref{sec:tabu} produced a valid
4-coloring of the Errera graph without backtracking. The central open question
is whether this extends to all planar graphs:

\begin{itemize}[noitemsep]
    \item Does the layer-by-layer tab\"{u} algorithm always produce a valid
        4-coloring for any planar graph, or are there counterexamples?
    \item If it fails on some graphs, can the within-layer ordering (currently
        ascending ID) be chosen more carefully --- for example by DSATUR within
        each layer --- to avoid failures?
    \item Is there a formal connection between the layer structure of the peeling
        and the reducible configurations of Robertson et al.~\cite{robertson1997}?
\end{itemize}

\subsection{Schnyder Woods and Grid Embedding}

\textbf{Schnyder woods}~\cite{schnyder1990} are an algebraic structure that
partitions the interior edges of a maximal triangulation into three spanning
trees $T_1, T_2, T_3$, one rooted at each outer-face vertex. At each interior
vertex $v$, exactly one outgoing edge of each color exists, and incoming and
outgoing edges alternate in a canonical cyclic order.

This structure generates integer barycentric coordinates: for each vertex $v$,
$a_i(v)$ counts faces in the region delimited by the $T_j$ and $T_k$ paths from
$v$ to their roots, with $a_1 + a_2 + a_3 = F - 1$.

\textbf{Coloring from Schnyder coordinates.} The parity map
$c(v) = (a_1(v) \bmod 2) \cdot 2 + (a_2(v) \bmod 2)$ gives a valid 4-coloring
of any maximal planar triangulation, provided the Schnyder wood satisfies the
exact local condition at every vertex. The key property is that for any edge
$\{u,v\}$ in a correct Schnyder wood, either $a_1(u) - a_1(v)$ or
$a_2(u) - a_2(v)$ is odd, guaranteeing $c(u) \neq c(v)$.

\textbf{Computational exploration.} We explored the connection between the
Schnyder coordinate ordering and the tab\"{u} algorithm. An approximate Schnyder
ordering was defined by sorting vertices by $(d_0+d_1+d_2, d_0, d_1)$ where
$d_i$ is the BFS distance from outer vertex $t_i$. This ordering approximates
the canonical ordering induced by the Schnyder wood. Combined with the recolor
strategy (Section~\ref{sec:forest}), this achieved \textbf{62.1\% 4-coloring}
on 1\,000 random Delaunay triangulations (vs 51\% for outer-face tab\"{u} with
recolor, and 21.5\% for plain outer-face tab\"{u}).

\textbf{Implementation status.} A full Schnyder wood requires the exact local
condition (cyclic pattern out-1/in-3/out-3/in-2/out-2/in-1 at every interior
vertex). Our BFS-approximate trees violate this condition for some edges, causing
the parity coloring to produce conflicts (adjacent vertices with equal
$(a_1 \bmod 2, a_2 \bmod 2)$). The correct implementation requires the canonical
ordering algorithm of de Fraysseix and Ossona de Mendez (1994), which carefully
tracks the ``active outer face'' at each step. This is left as future work.

\textbf{Open question.} Does a correct Schnyder ordering (processing vertices in
the order induced by the Schnyder wood's canonical sequence) as a tab\"{u}
selection rule, combined with the recolor strategy, approach or guarantee 100\%
4-coloring without backtracking? The approximate ordering already reduces
6-color failures from 9.2\% to 0.9\%, suggesting the exact ordering could be
significantly stronger.


\subsection{Discharging Method}

Implementing the discharging method would bridge the gap between the
constructive demonstration (backtracking produces a coloring) and the universal
proof (every planar graph admits one). The method assigns charge
$\mu(v) = 6 - \degr(v)$ to each vertex; the total sum is 12 by Euler's
formula. Discharge rules redistribute charge without changing the total; if no
reducible configuration is present in the graph, the redistribution produces a
contradiction ($\sum \mu'(v) \leq 0 < 12$)~\cite{robertson1997}.

\subsection{Reducibility Testing}

A configuration $C$ with outer ring $R$ is reducible if every valid 4-coloring
of $R$ can be extended to the interior of $C$. Verifying this requires
enumerating all colorings of the ring (at most $4^{|R|}$) and applying Kempe
swaps when necessary~\cite{robertson1997}. Implementing this test would connect
the system directly to the formal structure of Gonthier's proof.

% ══════════════════════════════════════════════════════════════════════════════
\section{Conclusions}
% ══════════════════════════════════════════════════════════════════════════════

We have documented a complete computational pipeline that, starting from the
representation of a planar graph as ordered pairs, produces a valid coloring
with at most four colors via:

\begin{enumerate}[noitemsep]
    \item Planarity verification in $O(|V|)$ (Boyer-Myrvold).
    \item Constrained maximalization preserving original edges.
    \item Tutte embedding: $14 \times 14$ Laplacian system for the Errera graph.
    \item Topological peeling by outer face, with elimination degrees $\leq 5$.
    \item Inverse greedy coloring guaranteeing at most 6 colors.
    \item Kempe chain analysis: identification of tangled chains.
    \item Layer-by-layer tab\"{u} coloring: 4 colors without backtracking on Errera.
    \item CSP backtracking (MRV + forward checking) with 4-color guarantee.
\end{enumerate}

The system concretely exhibits the failure mechanism of Kempe's algorithm on the
Errera graph (step 14: all chains tangled) and on the Kittell graph, which are
historically the first counterexamples to the incorrect proof of 1879. The gap
between the constructive coloring (backtracking produces the coloring) and the
universal proof (discharging proves one always exists) delimits the scope of the
system and points the direction of future work.

% ── Bibliography ──────────────────────────────────────────────────────────────
\newpage
\bibliographystyle{IEEEtran}
\begin{thebibliography}{99}

\bibitem{deita2024}
G. De Ita Luna and R. Marcial-Romero,
``A combinatorial algorithmic proof for the 4-coloring on planar graphs,''
Facultad de Ciencias de la Computaci\'on, BUAP, preprint, 2024.

\bibitem{appelhaken1977}
K. Appel and W. Haken,
``Every planar map is four colorable. Part I: Discharging,''
\textit{Illinois Journal of Mathematics}, vol. 21, no. 3, pp. 429--490, 1977.

\bibitem{boyermyrvold1999}
J. M. Boyer and W. J. Myrvold,
``On the cutting edge: Simplified $O(n)$ planarity by edge addition,''
\textit{Journal of Graph Algorithms and Applications}, vol. 8, no. 3,
pp. 241--273, 2004.

\bibitem{chrobak1995}
M. Chrobak and T. H. Payne,
``A linear-time algorithm for drawing a planar graph on a grid,''
\textit{Information Processing Letters}, vol. 54, no. 4, pp. 241--246, 1995.

\bibitem{chung1997}
F. R. K. Chung,
\textit{Spectral Graph Theory}.
Providence, RI: American Mathematical Society, 1997.

\bibitem{errera1921}
A. Errera,
``Du coloriage des cartes et de quelques questions d'analysis situs,''
Ph.D. dissertation, Universit\'e Libre de Bruxelles, 1921.

\bibitem{euler1752}
L. Euler,
``Elementa doctrinae solidorum,''
\textit{Novi Commentarii Academiae Scientiarum Imperialis Petropolitanae},
vol. 4, pp. 109--140, 1752.

\bibitem{fary1948}
I. F\'ary,
``On straight-line representation of planar graphs,''
\textit{Acta Scientiarum Mathematicarum}, vol. 11, pp. 229--233, 1948.

\bibitem{gonthier2008}
G. Gonthier,
``Formal proof---The four-color theorem,''
\textit{Notices of the American Mathematical Society}, vol. 55, no. 11,
pp. 1382--1393, 2008.

\bibitem{haralick1980}
R. M. Haralick and G. L. Elliott,
``Increasing tree search efficiency for constraint satisfaction problems,''
\textit{Artificial Intelligence}, vol. 14, no. 3, pp. 263--313, 1980.

\bibitem{heawood1890}
P. J. Heawood,
``Map-colour theorem,''
\textit{Quarterly Journal of Pure and Applied Mathematics}, vol. 24,
pp. 332--338, 1890.

\bibitem{kempe1879}
A. B. Kempe,
``On the geographical problem of the four colours,''
\textit{American Journal of Mathematics}, vol. 2, no. 3, pp. 193--200, 1879.

\bibitem{kittell1935}
I. Kittell,
``A group of operations on a partially colored map,''
\textit{Bulletin of the American Mathematical Society}, vol. 41, no. 6,
pp. 407--413, 1935.

\bibitem{kuratowski1930}
K. Kuratowski,
``Sur le probl\`eme des courbes gauches en topologie,''
\textit{Fundamenta Mathematicae}, vol. 15, pp. 271--283, 1930.

\bibitem{matula1983}
D. W. Matula and L. L. Beck,
``Smallest-last ordering and clustering and graph coloring algorithms,''
\textit{Journal of the ACM}, vol. 30, no. 3, pp. 417--427, 1983.

\bibitem{mohar2001}
B. Mohar and C. Thomassen,
\textit{Graphs on Surfaces}.
Baltimore, MD: Johns Hopkins University Press, 2001.

\bibitem{robertson1997}
N. Robertson, D. P. Sanders, P. Seymour, and R. Thomas,
``The four-colour theorem,''
\textit{Journal of Combinatorial Theory, Series B}, vol. 70, no. 1,
pp. 2--44, 1997.

\bibitem{russell2020}
S. Russell and P. Norvig,
\textit{Artificial Intelligence: A Modern Approach}, 4th ed.
Hoboken, NJ: Pearson, 2020.

\bibitem{schnyder1990}
W. Schnyder,
``Embedding planar graphs on the grid,''
in \textit{Proc. 1st ACM-SIAM Symposium on Discrete Algorithms (SODA)},
San Francisco, CA, 1990, pp. 138--148.

\bibitem{shamos1976}
M. I. Shamos and D. Hoey,
``Geometric intersection problems,''
in \textit{Proc. 17th Annual Symposium on Foundations of Computer Science (FOCS)},
Houston, TX, 1976, pp. 208--215.

\bibitem{tutte1963}
W. T. Tutte,
``How to draw a graph,''
\textit{Proceedings of the London Mathematical Society}, vol. 13, no. 1,
pp. 743--768, 1963.

\bibitem{wilson2002}
R. J. Wilson,
\textit{Four Colors Suffice: How the Map Problem Was Solved}.
Princeton, NJ: Princeton University Press, 2002.

\end{thebibliography}

\end{document}
